\scanpage{36}\mspage{35}
\deleted{Si on \deleted{\unsure{en|on}} a deux \unsure{comme ça|ensembles sc.} $F'$ et $F''$, correspondant à des ens d'indices $I',I''$, dire que $F''\preccurlyeq F'$ signifie que}

\deleted{\unsure{Donc}} Mais ds une telle situation,

Considérons l'ens ordonné suivant, défini en termes de la situation $(\deleted{\unsure{N}}*)$ ``avant $I'$''
\[
 K=J'\mathbin{\amalg} I\setminus J\qquad\text{en \deleted{\unsure{tant}} qu'ens.}
\]
la relation d'ordre sur $K$ induit celles données sur $J',I\setminus J$, et si $\alpha\in J'$, $\beta\in I-J$, on pose
\[
\left\{\begin{aligned}
 &\alpha\leq\beta\quad\text{sss}\quad h_J(\alpha)\leq\beta\quad\text{dans }I\\
 &\deleted{\unsure{\text{pour }\beta|\beta\leq\alpha}.}
\end{aligned}\right.
\]
donc $(\alpha,\beta)\in K\times K$, on a $\alpha\leq\beta$ sss
\margin{expliciter ds une proposition à part}
\[
\left|\begin{aligned}
 &\text{ou bien}\quad \deleted{\unsure{\alpha,\beta|\alpha}}\;(\alpha,\beta)\in\deleted{\unsure{I}}J' \ \text{et}\ \alpha\leq\beta\ \text{ds }J',\\
 &\text{ou bien}\quad (\alpha,\beta)\in I\setminus\deleted{\unsure{\alpha|\times|J}}J\ \text{et}\ \alpha\leq\beta\ \text{dans }I\setminus J\\
 &\text{ou bien}\quad \alpha\in\deleted{\unsure{I|J}}J',\ \beta\in I\setminus J\quad\text{et}\quad \deleted{\alpha\leq}h_J(\alpha)\leq\beta
\end{aligned}\right.
\]

Ceci posé, on a, en termes de $J'\xrightarrow{h_J}J\overset{i}{\hookrightarrow}I$ une factorisation canonique du composé
\diagram{\begin{tikzcd}[ampersand replacement=\&,column sep=large]
J'\arrow[r,hook,"{i}"]\arrow[d,"h_J"'] \& \overset{\deleted{\unsure{\text{iso}|I''|\text{en}}}}{K}\arrow[d,"h_K"]\\
J\arrow[r,hook,no head,"i"] \& I
\end{tikzcd}}
\margin{$h_K$ est $h_J$ sur $J'$ et l'\unsure{idel|ident.} can. sur $I\setminus J$}

On a fait ce qu'il fallait pour que sur $J'$, $K$ induise la relation de $J'$, que $h_K$ induise un iso d'ens ordonnés $K\setminus J'\xrightarrow{\sim}I\setminus J$, \deleted{et que de plus,} \inserted{cela implique que le} \deleted{\unsure{pour tout le} \unsure{G qui soit}} \inserted{carré est cartésien} ($J'=h_K^{-1}(J)$ avec ordre induit) D'autre part, on voit
